\documentclass[11pt]{article}

\usepackage{amsmath,amssymb,mathtools}
\usepackage[margin=1in]{geometry}

\newcommand{\RR}{\mathbb R}
\newcommand{\tridiag}{\operatorname{tridiag}}
\newcommand{\diag}{\operatorname{diag}}
\newcommand{\spec}{\operatorname{spec}}

\begin{document}

\begin{center}
{\Large Hatano--Nelson fourth-order recurrences on the real axis}
\end{center}

Fix
\[
0<a<1,\qquad
A_n(a)=\tridiag(a,0,1)\in\RR^{n\times n}.
\]
For real \(x\), set
\[
C_n(x):=xI_n-A_n(a),
\]
and for a squared singular-value parameter \(\lambda\in\RR\), define
\[
K_n(x,\lambda):=C_n(x)^TC_n(x)-\lambda I_n.
\]
Then
\[
s_n(x)^2=\lambda_{\min}\bigl(C_n(x)^TC_n(x)\bigr),
\]
and
\[
s_n(x)^2
\]
is the smallest root in \(\lambda\) of
\[
\det K_n(x,\lambda)=0.
\]

\section*{1. Pentadiagonal form}

Write
\[
c:=-(1+a)x,
\]
\[
B:=x^2+1+a^2-\lambda,
\]
\[
B_L:=x^2+a^2-\lambda,
\qquad
B_R:=x^2+1-\lambda.
\]
Then \(K_n(x,\lambda)\) is the symmetric pentadiagonal matrix with entries
\[
(K_n)_{j,j}=
\begin{cases}
B_L, & j=1,\\
B, & 2\le j\le n-1,\\
B_R, & j=n,
\end{cases}
\]
\[
(K_n)_{j,j+1}=(K_n)_{j+1,j}=c,
\]
and
\[
(K_n)_{j,j+2}=(K_n)_{j+2,j}=a.
\]

Thus, if
\[
K_n(x,\lambda)w=0,
\]
then the bulk equation is the fourth-order recurrence
\[
\boxed{
a w_{j-2}
+c w_{j-1}
+B w_j
+c w_{j+1}
+a w_{j+2}=0,
\qquad
3\le j\le n-2.
}
\tag{1}
\]
The boundary equations are
\[
\boxed{
B_Lw_1+cw_2+aw_3=0,
}
\tag{2}
\]
\[
\boxed{
cw_1+Bw_2+cw_3+aw_4=0,
}
\tag{3}
\]
\[
\boxed{
aw_{n-3}+cw_{n-2}+Bw_{n-1}+cw_n=0,
}
\tag{4}
\]
and
\[
\boxed{
aw_{n-2}+cw_{n-1}+B_Rw_n=0.
}
\tag{5}
\]
Terms with indices outside \(\{1,\dots,n\}\) are omitted in the evident small-\(n\) cases.

\section*{2. Bulk factorization}

Let
\[
(\mathcal L w)_j:=w_{j+1}+w_{j-1}.
\]
Then the bulk operator in (1) factors as
\[
\boxed{
a(\mathcal L-y_+)(\mathcal L-y_-)w=0,
}
\tag{6}
\]
where \(y_\pm\) are the two roots of
\[
a y^2-(1+a)xy+x^2+(1-a)^2-\lambda=0.
\]
Explicitly,
\[
\boxed{
y_\pm=
\frac{(1+a)x
\pm
\sqrt{(1-a)^2(x^2-4a)+4a\lambda}}
{2a}.
}
\tag{7}
\]
Equivalently,
\[
\boxed{
y_++y_-=\frac{(1+a)x}{a},
}
\tag{8}
\]
and
\[
\boxed{
y_+y_-=
\frac{x^2+(1-a)^2-\lambda}{a}.
}
\tag{9}
\]

Let
\[
P_j(y):=U_j(y/2),
\]
where \(U_j\) is the Chebyshev polynomial of the second kind. Thus
\[
P_{j+1}(y)=yP_j(y)-P_{j-1}(y),
\qquad
P_{-1}(y)=0,\quad P_0(y)=1.
\]
The two second-order channels are generated by
\[
P_j(y_+),\qquad P_j(y_-).
\]

\section*{3. Product-channel recurrence}

Define
\[
S_j:=P_j(y_+)P_j(y_-).
\]
Let
\[
p:=y_+y_-,
\]
and
\[
q:=y_+^2+y_+y_-+y_-^2-2.
\]
In terms of \(a,x,\lambda\), these are
\[
\boxed{
p=
\frac{x^2+(1-a)^2-\lambda}{a},
}
\tag{10}
\]
and
\[
\boxed{
q=
\frac{(1+a+a^2)x^2+a\lambda-a(1+a^2)}{a^2}.
}
\tag{11}
\]
The generating function of \(S_j\) is
\[
\boxed{
S(z):=\sum_{j\ge0}S_jz^j
=
\frac{1-z^2}{\chi(z)},
}
\tag{12}
\]
where
\[
\boxed{
\chi(z)
=
1-pz+(q-p)z^2-pz^3+z^4.
}
\tag{13}
\]
Therefore \(S_j\) satisfies the fourth-order palindromic recurrence
\[
\boxed{
S_j
-pS_{j-1}
+(q-p)S_{j-2}
-pS_{j-3}
+S_{j-4}=0,
\qquad j\ge4.
}
\tag{14}
\]

\section*{4. Determinant formula}

Normalize the determinant by
\[
\det K_n(x,\lambda)=a^{\,n-1}\mathcal E_n(x,\lambda).
\]
Then
\[
\boxed{
\mathcal E_n
=
aS_n
-(1-a)^2
\sum_{j=0}^{n-1}(n-j)S_j.
}
\tag{15}
\]
Consequently,
\[
\boxed{
s_n(x)^2
=
\min\{\lambda\in\RR:\mathcal E_n(x,\lambda)=0\}.
}
\tag{16}
\]

Equivalently, \(\mathcal E_n\) is given by the rational generating function
\[
\boxed{
\sum_{n\ge0}\mathcal E_nz^n
=
\frac{(1+z)(a-z)(1-az)}
{(1-z)\chi(z)}.
}
\tag{17}
\]
Let
\[
\mathfrak D(z):=(1-z)\chi(z).
\]
Then
\[
\boxed{
\mathfrak D(z)
=
1-(p+1)z+qz^2-qz^3+(p+1)z^4-z^5.
}
\tag{18}
\]
Thus \(\mathcal E_n\) satisfies the fifth-order anti-palindromic recurrence
\[
\boxed{
\mathcal E_n
-(p+1)\mathcal E_{n-1}
+q\mathcal E_{n-2}
-q\mathcal E_{n-3}
+(p+1)\mathcal E_{n-4}
-\mathcal E_{n-5}
=0,
}
\tag{19}
\]
for all \(n\) beyond the initial numerator range.

The first values are
\[
\boxed{
\mathcal E_0=a,
}
\tag{20}
\]
\[
\boxed{
\mathcal E_1=x^2-\lambda,
}
\tag{21}
\]
and
\[
\boxed{
\mathcal E_2
=
\frac{
(x^2+a^2-\lambda)(x^2+1-\lambda)
-(1+a)^2x^2
}{a}.
}
\tag{22}
\]

\section*{5. Auxiliary coefficient recurrence}

Define
\[
J(z):=
\frac{(a-z)(1-az)}{\mathfrak D(z)}
=
\sum_{m\ge0}J_mz^m.
\]
Then
\[
\boxed{
\mathcal E_n=J_n+J_{n-1},
\qquad J_{-1}:=0.
}
\tag{23}
\]
The coefficients \(J_m\) satisfy
\[
\boxed{
J_m
-(p+1)J_{m-1}
+qJ_{m-2}
-qJ_{m-3}
+(p+1)J_{m-4}
-J_{m-5}
=r_m,
}
\tag{24}
\]
where
\[
r_0=a,\qquad
r_1=-(1+a^2),\qquad
r_2=a,\qquad
r_m=0\quad(m\ge3),
\]
and \(J_m=0\) for \(m<0\).

\section*{6. Symmetric banded \(LDL^T\) recurrence}

The smallest squared singular value can also be tracked by symmetric banded elimination.

Let
\[
b_1:=B_L=x^2+a^2-\lambda,
\]
\[
b_j:=B=x^2+1+a^2-\lambda,
\qquad 2\le j\le n-1,
\]
and
\[
b_n:=B_R=x^2+1-\lambda.
\]
Set
\[
c:=-(1+a)x.
\]
Define the pivots \(\alpha_j\), first subdiagonal multipliers \(\ell_j\), and second subdiagonal multipliers \(m_j\) by
\[
\boxed{
\alpha_1=b_1,
}
\tag{25}
\]
\[
\boxed{
\ell_1=\frac{c}{\alpha_1},
\qquad
m_1=\frac{a}{\alpha_1}.
}
\tag{26}
\]
For \(j\ge2\),
\[
\boxed{
\alpha_j
=
b_j-\ell_{j-1}^2\alpha_{j-1}-m_{j-2}^2\alpha_{j-2},
}
\tag{27}
\]
with the convention \(m_0=0\). Also,
\[
\boxed{
\ell_j
=
\frac{c-a\ell_{j-1}}{\alpha_j},
}
\tag{28}
\]
and
\[
\boxed{
m_j=\frac{a}{\alpha_j}.
}
\tag{29}
\]
Then
\[
\boxed{
\det K_n(x,\lambda)=\prod_{j=1}^n\alpha_j.
}
\tag{30}
\]
Moreover,
\[
\boxed{
K_n(x,\lambda)\succ0
\quad\Longleftrightarrow\quad
\alpha_1,\dots,\alpha_n>0.
}
\tag{31}
\]
Hence
\[
\boxed{
s_n(x)^2
=
\sup\{\lambda:\alpha_1(x,\lambda),\dots,\alpha_n(x,\lambda)>0\}.
}
\tag{32}
\]

\section*{7. Repaired right-end recurrence}

Define the repaired matrix
\[
H_n(x,\lambda):=K_n(x,\lambda)+a^2e_ne_n^T.
\]
Then the right terminal diagonal of \(H_n\) is repaired from
\[
B_R=x^2+1-\lambda
\]
to
\[
B=x^2+1+a^2-\lambda.
\]
Thus the repaired terminal row satisfies
\[
\boxed{
aw_{n-2}+cw_{n-1}+Bw_n=0.
}
\tag{33}
\]
The unrepaired old matrix is the rank-one downdate
\[
\boxed{
K_n=H_n-a^2e_ne_n^T.
}
\tag{34}
\]

For the right append \(n\to n+1\), the leading block of \(K_{n+1}\) is \(H_n\), and
\[
\boxed{
K_{n+1}
=
\begin{pmatrix}
H_n & b\\
b^T & d_R
\end{pmatrix},
}
\tag{35}
\]
where
\[
\boxed{
b=a e_{n-1}+c e_n
=
a e_{n-1}-(1+a)x e_n,
}
\tag{36}
\]
and
\[
\boxed{
d_R=x^2+1-\lambda.
}
\tag{37}
\]

\section*{8. Two-site Schur tail recurrence}

Eliminate the repaired prefix
\[
P:=H_{n-1}.
\]
Write
\[
q:=a e_{n-2}+c e_{n-1},
\qquad
f:=a e_{n-1},
\]
\[
h:=B=x^2+1+a^2-\lambda,
\qquad
d:=B_R=x^2+1-\lambda.
\]
Then the final two-site Schur tail is
\[
\boxed{
T=
\begin{pmatrix}
\sigma & \chi\\
\chi & D
\end{pmatrix},
}
\tag{38}
\]
where
\[
\boxed{
\sigma=h-q^TP^{-1}q,
}
\tag{39}
\]
\[
\boxed{
\chi=c-q^TP^{-1}f,
}
\tag{40}
\]
and
\[
\boxed{
D=d-f^TP^{-1}f.
}
\tag{41}
\]
The repaired old block satisfies
\[
H_n\succ0
\quad\Longleftrightarrow\quad
P\succ0,\quad \sigma>0.
\]
The unrepaired old block satisfies
\[
K_n\succ0
\quad\Longleftrightarrow\quad
P\succ0,\quad \sigma>a^2.
\]
The appended block satisfies
\[
K_{n+1}\succ0
\quad\Longleftrightarrow\quad
P\succ0,\quad
\sigma>0,\quad
\sigma D-\chi^2>0.
\]

\section*{9. Summary of the main fourth-order objects}

The Hatano--Nelson real-axis singular-value problem is governed by the fourth-order bulk recurrence
\[
a w_{j-2}
-(1+a)xw_{j-1}
+
(x^2+1+a^2-\lambda)w_j
-(1+a)xw_{j+1}
+
a w_{j+2}=0.
\]
It factors as
\[
a(\mathcal L-y_+)(\mathcal L-y_-)w=0,
\]
with
\[
y_\pm=
\frac{(1+a)x
\pm
\sqrt{(1-a)^2(x^2-4a)+4a\lambda}}
{2a}.
\]
The determinant is
\[
\det K_n(x,\lambda)=a^{n-1}\mathcal E_n(x,\lambda),
\]
where
\[
\sum_{n\ge0}\mathcal E_nz^n
=
\frac{(1+z)(a-z)(1-az)}
{
1-(p+1)z+qz^2-qz^3+(p+1)z^4-z^5
}.
\]
Finally,
\[
s_n(x)^2
\]
is the smallest root in \(\lambda\) of
\[
\mathcal E_n(x,\lambda)=0.
\]

\end{document}